\par\Needspace{6\baselineskip}\medskip
\hypertarget{native-125-article092}{}\pdfbookmark[2]{92.}{article-092}
\begin{center}92.\end{center}
Although equations 15 and $15^{\ast}$ can have three real roots in certain cases, there will never be any doubt as to which one to take in our problem. Since $h$ is clearly positive, it follows readily from the theory of equations that equation 15 has just one positive root together with either two imaginary or two negative roots. Since $y=\dfrac{m}{\BareRoot{(l+x)}}$ must be positive, clearly no uncertainty \SourceLeaf{126}remains here. As for equation $15^{\ast}$, we first note that $L$ must be greater than 1. This is easily proved by writing the equation given in Art. 89 in the form $L=1+\dfrac{\cos\tfrac12 f^2}{-\cos f}+\dfrac{\operatorname{tang}2\omega^2}{-\cos f}$. Further, substituting in equation $12^{\ast}$ for $M$ the expression $Y\BareRoot{(L-x)}$ gives $Y+1=(L-x)X$, and hence
\[
\begin{aligned}
Y+1>(1-x)X>{}&\tfrac43+\dfrac{4}{3\mathbin{.}5}x+\dfrac{4\mathbin{.}6}{3\mathbin{.}5\mathbin{.}7}xx
\\[4pt]
&+\dfrac{4\mathbin{.}6\mathbin{.}8}{3\mathbin{.}5\mathbin{.}7\mathbin{.}9}x^3+\text{etc.}>\tfrac43,
\end{aligned}
\]
and consequently $Y>\tfrac13$. Thus, if we set $Y=\tfrac13+Y'$, then $Y'$ must be positive, and equation $15^{\ast}$ becomes
\[
Y'^3+2Y'Y'+(1-H)Y'+\tfrac{4}{27}-\tfrac29H=0,
\]
which, as the theory of equations readily shows, cannot have more than one positive root. It follows that equation $15^{\ast}$ has exactly one root greater than $\tfrac13$\footnote{Provided that we assume the problem really admits a solution.}, and this is the only root we take for our problem\nobreak\HistoricalNoteDouble{Manuscript note:] Equation $15^{\ast}$ has
\par 1) One real negative root, along with two imaginary roots, when $H$ lies between the limits 0 and $\dfrac{+1+\BareRoot{5}}{6}$.
\par 2) Three real roots, one of them positive, when $H$ lies between $\dfrac{+1-\BareRoot{5}}{6}$ and 0.
\par 3) One real positive root and two imaginary roots when $H$ lies between $-\infty$ and $\dfrac{+1-\BareRoot{5}}{6}\mathpunct{\cdot}$
\par 4) Three real roots, one of them negative, when $H$ lies between $\dfrac{+1+\BareRoot{5}}{6}$ and $+\infty$.
\par Clearly, then, only case 4 can apply here, in which there are two positive roots. One of them is always less than $\dfrac{-1+\BareRoot{5}}{6}\mathpunct{\raisebox{.45ex}{,}}$ and the other is greater, so only the latter can be admissible.
\par From the procedure indicated in the text, it is clear that the last equation could have two positive roots only if $1-H$ were negative and $\tfrac{4}{27}-\tfrac29H$ positive at the same time. This is clearly impossible, since $\tfrac{4}{27}-\tfrac29H=\tfrac29(1-H)-\tfrac{2}{27}$.
}.
